% Focused framework presentation.
\section[Research directions]{Algebraic tools and subsequent research}

\begin{frame}[t]{Algebraic tools supplied by the framework}
\hypertarget{main:tools}{}
\heading{A trajectory equation becomes a kernel and a translate}
\begin{equation*}
  \mathsf D_A\gamma=B\quad\Longrightarrow\quad
  \mathcal S_B=\gamma_{\mathrm p}+\ker\mathsf D_A.
\end{equation*}
\heading{Evaluation turns uniqueness into composition}
\begin{equation*}
  \Phi_{s,t}=E_tE_s^{-1},\qquad
  \Phi_{r,t}\Phi_{s,r}=E_tE_r^{-1}E_rE_s^{-1}.
\end{equation*}
\heading{Pullback turns state motion into a linear action}
\begin{equation*}
  \mathcal K_{s,t}f=f\circ\Phi_{s,t},\qquad
  \langle\mathcal K_{s,t}f,\mu\rangle=\langle f,(\Phi_{s,t})_\#\mu\rangle.
\end{equation*}
\heading{Trace, determinant, and tilting simplify derived problems}
The trace reduces local volume evolution to a scalar ODE.
Gibbs tilting reduces a measure optimization problem to KL nonnegativity.
\slidecite{Synthesis of the constructions in manuscript \S\S\,1.2.1--1.2.4.}
\end{frame}

\begin{frame}[t]{Research direction: projection to collective variables}
\hypertarget{main:research-projection}{}
Choose $\xi:V\to\mathcal Z$ and write $Z_t=\xi(X_t)$, $\nu_t=\xi_\#\mu_t$.
\heading{Pullback embeds reduced observables}
\begin{equation*}
  C_\xi g=g\circ\xi.
\end{equation*}
On $L^2$, conditional expectation is an orthogonal projection onto functions of $Z_t$.
\heading{A candidate reduced generator}
For suitable test functions and integrable terms, define
\begin{equation*}
  (\overline{\mathcal L}_t g)(z)
  =\mathbb E[(\mathcal L_t C_\xi g)(X_t)\mid Z_t=z].
\end{equation*}
This conditional projection matches the instantaneous evolution of reduced expectations.
\heading{Research question}
Determine when it defines a well-posed reduced process, and measure the information
lost in temporal correlations. Matching one-time laws alone does not establish dynamical closure.
\slidecite{Research direction based on manuscript \S\,1.3, collective variables and conditional effective dynamics.}
\end{frame}

\begin{frame}[t]{Research direction: conditional learning of flows}
\hypertarget{main:research-learning}{}
For a differentiable interpolation $I_t=\alpha_tX_0+\lambda_tX_1$, seek a velocity
that transports its one-time laws.
\heading{A projection supplies the population target}
If $\dot I_t$ is square integrable, then
\begin{equation*}
  \begin{aligned}
    u_t^*&=\operatorname*{arg\,min}_{u\in L^2(\operatorname{Law}(I_t))}
      \mathbb E\|u(I_t)-\dot I_t\|^2,\\
    u_t^*(x)&=\mathbb E[\dot I_t\mid I_t=x].
  \end{aligned}
\end{equation*}
The conditional projection defines the target almost everywhere.
\heading{Proposed extension}
Fit a velocity or Hamiltonian in a class encoding the required structure.
Compare induced operators and transported laws with the analytic baselines.
\heading{What must be checked}
Separate representation, sampling, optimization, and time-discretization errors.
The population identity does not guarantee exact recovery by a finite model.
\slidecite{Research direction: manuscript \S\,1.4, conditional learning and structural approximation.}
\end{frame}

\begin{frame}[t]{A research programme built on the framework}
\hypertarget{main:research-programme}{}
\heading{Mathematical foundation}
Use the same state--observable--law pairing to formulate exact evolution,
reduction, learning, and numerical approximation.
\heading{Next validation tasks}
\begin{enumerate}\setlength{\itemsep}{8pt}
  \item Use affine flows and the coupled OU benchmark to compare exact and reduced generators.
  \item Test both transported one-time laws and two-time correlations when assessing a learned reduction.
  \item Check the relevant invariants under discretization: probability mass and positivity, or Hamiltonian symplecticity.
\end{enumerate}
\heading{The organizing question}
Which algebraic identities survive each approximation, and what error appears when they fail?
\par\medskip
The proposed extensions use the exact constructions as references for measuring
the errors introduced by reduction and learning.
\slidecite{Research programme based on manuscript \S\S\,1.3--1.5 and the existing benchmark design.}
\end{frame}
